% Caching, Performance, and Optimization

\section{Redis Caching Strategy}

\subsection{Cache Implementation with TTL and LRU}

\begin{lstlisting}[language=JavaScript, caption=Redis Cache with Expiration]
const redis = require('redis');

const cacheClient = redis.createClient({
  host: 'localhost',
  port: 6379,
  retry_strategy: (opts) => {
    if (opts.total_retry_time > 1000 * 60 * 60) {
      return new Error('Retry time exhausted');
    }
    return Math.min(opts.attempt * 100, 3000);
  }
});

/**
 * Generic cache wrapper
 * @param {string} key - Cache key
 * @param {function} fetchFunction - Function to fetch data if not cached
 * @param {number} ttl - Time to live in seconds (default: 3600)
 */
const getCachedData = async (key, fetchFunction, ttl = 3600) => {
  try {
    // Check cache first
    const cached = await cacheClient.get(key);
    if (cached) {
      console.log(`✓ Cache HIT for key: ${key}`);
      return JSON.parse(cached);
    }

    console.log(`✗ Cache MISS for key: ${key}`);
    
    // Fetch fresh data
    const data = await fetchFunction();
    
    // Store in cache with TTL
    await cacheClient.setex(
      key,
      ttl,
      JSON.stringify(data)
    );
    
    return data;
  } catch (error) {
    console.error(`Cache error for key ${key}:`, error);
    throw error;
  }
};

// Usage example: Cache API responses
const getUser = async (userId) => {
  return getCachedData(
    `user:${userId}`,
    async () => {
      const response = await fetch(\`/api/users/\${userId}\`);
      return response.json();
    },
    3600  // Cache for 1 hour
  );
};
\end{lstlisting}

\subsection{Cache Invalidation Strategies}

\begin{lstlisting}[language=JavaScript, caption=Cache Invalidation Patterns]
// Pattern 1: Manual Invalidation
const invalidateCache = async (key) => {
  await cacheClient.del(key);
  console.log(`Cache invalidated: ${key}`);
};

// Pattern 2: Cascade Invalidation
const invalidateCascade = async (pattern) => {
  const keys = await cacheClient.keys(pattern);
  if (keys.length > 0) {
    await cacheClient.del(...keys);
    console.log(`Invalidated ${keys.length} keys matching: ${pattern}`);
  }
};

// Pattern 3: Event-based Invalidation
const handleUserUpdate = async (userId) => {
  await invalidateCascade(`user:${userId}:*`);
  await invalidateCascade(`users:list:*`);
};

// Pattern 4: Conditional Refresh
const refreshIfStale = async (key, fetchFunction, maxAge = 300) => {
  const ttl = await cacheClient.ttl(key);
  
  if (ttl <= 0) {
    // Missing or expired
    return getCachedData(key, fetchFunction);
  } else if (ttl < maxAge) {
    // Refresh in background
    getCachedData(key, fetchFunction).catch(err => 
      console.error('Background refresh failed:', err)
    );
  }
  
  const cached = await cacheClient.get(key);
  return JSON.parse(cached);
};
\end{lstlisting}

\section{Rate Limiting and API Protection}

\subsection{Express Rate Limiting Middleware}

\begin{lstlisting}[language=JavaScript, caption=API Rate Limiting]
const rateLimit = require('express-rate-limit');
const RedisStore = require('rate-limit-redis');

// Create Redis store for distributed rate limiting
const redisStore = new RedisStore({
  client: cacheClient,
  prefix: 'rl:'
});

// General API limiter
const apiLimiter = rateLimit({
  store: redisStore,
  windowMs: 15 * 60 * 1000,  // 15 minutes
  max: 100,  // Limit each IP to 100 requests per windowMs
  message: 'Too many requests from this IP',
  standardHeaders: true,  // Send rate limit info in headers
  legacyHeaders: false
});

// Strict limiter for auth endpoints
const authLimiter = rateLimit({
  store: redisStore,
  windowMs: 15 * 60 * 1000,
  max: 5,  // 5 attempts
  skipSuccessfulRequests: true  // Don't count successful requests
});

// Apply limiters to routes
app.use('/api/', apiLimiter);
app.post('/auth/login', authLimiter, loginController);
app.post('/auth/register', authLimiter, registerController);
\end{lstlisting}

\subsection{Request Throttling for Email Service}

\begin{lstlisting}[language=JavaScript, caption=Email Throttling]
const pThrottle = require('p-throttle');
const nodemailer = require('nodemailer');

// Throttle to 1 email per second
const emailThrottle = pThrottle({
  limit: 1,
  interval: 1000
});

const mailer = nodemailer.createTransporter({
  service: 'gmail',
  auth: {
    user: process.env.EMAIL_USER,
    pass: process.env.EMAIL_PASS
  }
});

const sendThrottledEmail = emailThrottle(async (to, subject, html) => {
  try {
    const info = await mailer.sendMail({
      from: process.env.EMAIL_USER,
      to,
      subject,
      html
    });
    console.log('Email sent:', info.messageId);
    return info;
  } catch (error) {
    console.error('Failed to send email:', error);
    throw error;
  }
});

// Usage: Send bulk emails safely
const sendBulkEmails = async (recipients) => {
  const promises = recipients.map(recipient =>
    sendThrottledEmail(
      recipient.email,
      'Newsletter',
      '<h1>Welcome</h1>'
    )
  );
  return Promise.all(promises);
};
\end{lstlisting}

\newpage
